%=========== ΛΥΣΗ - 38ο ΘΕΜΑ Α ===========
\providecommand{\theoriaref}[3]{%
\begin{center}
\fcolorbox{maincolor!70!black}{maincolor!8}{%
\begin{minipage}{.88\linewidth}
\textcolor{maincolor!80!black}{\kerkissans{\textbf{#1~\ref{#2}}}}\\[1mm]
#3
\end{minipage}}
\end{center}}
\providecommand{\orismosref}[1]{%
\theoriaref{Ορισμός}{#1}{Η απάντηση δίνεται από τον αντίστοιχο ορισμό.}}
\providecommand{\apodeixiref}[1]{%
\theoriaref{Θεώρημα}{#1}{Η ζητούμενη απόδειξη δίνεται στην αντίστοιχη απόδειξη.}}

\begin{thema}{Α}
\begin{erwthma}
\item \apodeixiref{thewrhma:paragwgos_dynamhs_arnhtikos}

\item \orismosref{orismos:synarthsh_1_1}

\item Οι ζητούμενοι ορισμοί είναι:
\orismosref{orismos:oliko_megisto}
\orismosref{orismos:oliko_elaxisto}

\item Οι χαρακτηρισμοί των προτάσεων είναι:
\begin{alist}
\item \textbf{Σωστό}. Αν στο σημείο καμπής η συνάρτηση είναι παραγωγίσιμη, τότε η εφαπτομένη στο σημείο αυτό διαπερνά την καμπύλη.
\item \textbf{Λάθος}. Σε ένωση ξένων διαστημάτων δύο παράγουσες μπορεί να διαφέρουν κατά διαφορετική σταθερά σε κάθε διάστημα.
\item \textbf{Λάθος}. Η γραφική παράσταση της $f^{-1}$ είναι συμμετρική της $C_f$ ως προς την ευθεία $y=x$, όχι ως προς τον άξονα $x'x$.
\item \textbf{Σωστό}. Για $f(x)=\ln x$ ισχύει
\[ \lim_{x\to0^+}\ln x=-\infty, \]
άρα ο άξονας $y'y$, δηλαδή η ευθεία $x=0$, είναι κατακόρυφη ασύμπτωτη.
\item \textbf{Σωστό}. Έχουμε
\[ \dint_1^2 e^x\d x=e^2-e. \]
\end{alist}

\item Σωστή απάντηση είναι η \textbf{β}. Θέτοντας $u=2x$, έχουμε $u\to0$ και
\[ \dfrac{f(2x)}{x}=2\cdot\dfrac{f(2x)}{2x}=2\cdot\dfrac{f(u)}{u}\to 2\cdot2=4. \]
\end{erwthma}
\end{thema}
%=========== ΤΕΛΟΣ ΛΥΣΗΣ - 38ο ΘΕΜΑ Α ===========
